problem:  
Let \( (M,g) \) be a closed Riemannian surface and \( (N,h) \) a compact Riemannian manifold.  
For each \(n\), let \(G_n=(V_n,E_n,F_n)\) be a geodesic triangulation of \(M\) with mesh size \(h_n\to 0\).  
Define a **discrete generalized harmonic map** \(f_n:V_n\to N\) as a minimizer of the discrete energy  
\[
E_n(f_n)=\frac12\!\!\sum_{(i,j)\in E_n} \!w_{ij}^{(n)}\, d_N^2(f_n(i),f_n(j)),
\]
where the symmetric weights \(w_{ij}^{(n)}>0\) are chosen so that \(E_n\) \(\Gamma\)-converges to the Dirichlet energy of maps \(f:M\to N\).  
Assume \(f_n\) converges (weakly in \(H^1\)) to a smooth harmonic map \(f:M\to N\).

Denote by
\[
\tau_{ij}^{(n)} = w_{ij}^{(n)}\, \exp^{-1}_{f_n(i)}(f_n(j)) \in T_{f_n(i)}N
\]
the *discrete tension vectors* along edges of \(G_n\).  

**Open problem:**  
Does there exist a coordinate-free, mesh-invariant functional
\[
\Phi\big(\{\tau^{(n)}_{ij}\}_{(i,j)\subset \partial T}\big)
\]
defined per triangle \(T \in F_n\), depending *only* on the local discrete tension field and weights (not on the full metric geometry of \(N\)), such that the resulting **discrete curvature density**
\[
\kappa_n(T) = \Phi\big(\{\tau^{(n)}_{ij}\}_{(i,j)\subset \partial T}\big)
\]
satisfies the curvature convergence property
\[
\kappa_n(T) \xrightarrow[n\to\infty]{} K_N(df_p(e_1)\wedge df_p(e_2))
\]
for almost every \(p\in M\) associated with the limiting location of the triangles \(T\), where \(K_N\) is the sectional curvature of \(N\) and \(df_p(e_1),df_p(e_2)\) span the image plane of \(df_p\)?  

Equivalently, can intrinsic curvature information of the target manifold \(N\) be recovered asymptotically *purely from the local discrete harmonic energy structure*—without directly sampling metric distances in \(N\)?

Why is it a "good" problem:  
This question focuses sharply on the interplay between **first-order variational structure** and **second-order curvature geometry**. It asks whether curvature can emerge as a universal limit property encoded implicitly in the equilibrium data of discrete harmonic maps, rather than as an explicit geometric measurement.  

The formulation is precise (defining triangulations, convergence, and the curvature target) and conceptually deep: it links discrete harmonic analysis, Γ-convergence, and intrinsic curvature reconstruction. It also imposes a crucial restriction—Φ must depend only on discrete tension data—to ensure the result cannot be reduced to trivial metric-triangle curvature computations.  

A positive resolution would reveal that purely variational energy structures are sufficient to reconstruct curvature, unifying discrete geometric analysis and Riemannian geometry at a fundamental level. A negative answer would delineate intrinsic limits of variational discretizations. Either outcome would profoundly clarify the relationship between discrete harmonic theory and continuous curvature, making this a “narrow entrance, profound depth” kind of problem and thus an exemplary "good" mathematical problem.